\section*{Chapitre \ref{ch:iv1:timed-test-4} --- Épreuve chronométrée 4}
\markboth{\protect\omTocMark{}{Réponses~: Épreuve chronométrée 4}}{\protect\omTocMark{}{Réponses~: Épreuve chronométrée 4}}

\begin{solution}{iq:iv1:timed-test-4:1}
Aucune correspondance n'a une probabilité $\frac{8}{21}$, et exactement une paire a une probabilité $\frac47$.

En tirant les chaussettes une à une, la première peut être quelconque. La deuxième doit éviter son partenaire~: 8 choix sur 9 restent possibles. La troisième doit éviter deux partenaires~: 6 choix sur 8. La quatrième doit en éviter trois~: 4 choix sur 7. Le produit vaut $\frac89 \cdot \frac68 \cdot \frac47 = \frac{8}{21}$\index{comptage séquentiel}.

Exactement une paire~: on choisit la paire, de 5 façons, puis 2 des 4 autres paires, et une chaussette de chacune~: $5 \times \binom42 \times 4 = 120$, sur $\binom{10}{4} = 210$. On obtient $\frac47$.

Deux paires se présentent de $\binom52 = 10$ façons, soit $\frac{1}{21}$, et les trois cas somment à 1, ce qui constitue une vérification utile.

\iqwrong{Compter les tirages ordonnés à un endroit et les tirages non ordonnés à un autre, ce qui est l'erreur habituelle sur le second résultat.}
\end{solution}

\begin{solution}{iq:iv1:timed-test-4:2}
La probabilité vaut $\frac14$.

Sachant $Y + Z = s$, on a $P(X > s) = e^{-s}$. La réponse est donc $\mathbb{E}e^{-(Y + Z)} = (\mathbb{E}e^{-Y})^2$. Et $\mathbb{E}e^{-Y} = \int_0^\infty e^{-2y}\,dy = \frac12$. Le résultat vaut $\frac14$\index{absence de mémoire}.

Une seconde méthode passe par une course. Considérons trois horloges exponentielles. Le premier signal parmi $X$ et $Y$ provient de $Y$ avec probabilité $\frac12$. Par la propriété d'absence de mémoire, $X$ repart à zéro. Il faut alors que $Z$ batte aussi ce temps, encore avec probabilité $\frac12$. D'où $\frac12 \times \frac12$. Une simulation le confirme.

\iqfollow{Quelle est la valeur de $P(X > 2Y)$~?}{$\mathbb{E}e^{-2Y} = \int_0^\infty e^{-3y}\,dy = \frac13$.}
\end{solution}

\enlargethispage{-3\baselineskip}%
\begin{solution}{iq:iv1:timed-test-4:3}
La pente vaut $\frac12$, l'ordonnée à l'origine vaut $\frac16$, et $R^2 = \frac34$.

Les moyennes valent $\bar{x} = 1$ et $\bar{y} = \frac23$. On a alors $S_{xx} = 1 + 0 + 1 = 2$ et $S_{xy} = (-1)(-\frac23) + 0 + (1)(\frac13) = 1$. La pente vaut donc $\frac12$\index{moindres carrés}. La droite passe par le point moyen, d'où une ordonnée à l'origine de $\frac23 - \frac12 = \frac16$.

La somme totale des carrés vaut $\frac49 + \frac19 + \frac19 = \frac23$, et la part expliquée vaut $b^2S_{xx} = \frac12$. D'où $R^2 = \frac12/\frac23 = \frac34$.

Les résidus valent $-\frac16$, $\frac13$ et $-\frac16$. Leur somme est nulle, et ils sont orthogonaux à $x$, comme il se doit. Cela permet de repérer une erreur de calcul.

\iqwrong{Imposer à la droite de passer par l'origine, ce qui donne une pente de $\frac35$. Rien dans l'énoncé n'autorise à supprimer l'ordonnée à l'origine.}
\end{solution}

\begin{solution}{iq:iv1:timed-test-4:4}
Le nombre espéré vaut 6, quelle que soit la case de départ située en dessous de 6.

Notons $e_k$ le nombre espéré de lancers depuis la case $k$. Depuis 5, seul un 1 termine la partie, donc $e_5 = 6$.

Depuis 4~: un 1 conduit à 5, un 2 termine la partie, et les quatre autres lancers laissent sur place. On a donc $e_4 = 1 + \frac16e_5 + \frac46e_4$, ce qui donne $e_4 = 6$. Par récurrence\index{analyse au premier pas}, si $e_{k + 1}, \dots, e_5$ valent tous 6, alors depuis $k$ il reste $6 - k$ lancers utiles, dont un termine la partie et les autres mènent à des cases de valeur 6. Il vient
\[
e_k = 1 + \frac{5 - k}{6}6 + \frac{k}{6}e_k,
\]
ce qui donne $e_k = 6$.

La raison est simple~: depuis n'importe quelle case, exactement un lancer termine la partie à ce tour. Chaque lancer réussit donc avec probabilité $\frac16$, quelle que soit la case, et l'attente suit une loi géométrique. Une simulation le confirme.

\iqwrong{Calculer une longue récurrence sans jamais remarquer le motif. Une réponse surprenante de 6 doit inciter à chercher la raison simple.}
\end{solution}

\begin{solution}{iq:iv1:timed-test-4:5}
La distance moyenne vaut $\frac{4}{\pi} \approx 1.27$.

Fixons le premier point. L'angle $\theta$ vers le second suit une loi uniforme sur $[0, 2\pi]$. La corde mesure $2\sin(\theta/2)$\index{longueur de corde}. D'où
\[
\mathbb{E} = \frac{1}{2\pi}\int_0^{2\pi}2\sin\frac{\theta}{2}\,d\theta = \frac{1}{2\pi}\left[-4\cos\frac{\theta}{2}\right]_0^{2\pi} = \frac{8}{2\pi}.
\]
Un contrôle de cohérence~: la distance est comprise entre 0 et 2, et une corde à angle droit mesure $\sqrt{2} = 1.41$. La moyenne est inférieure à cette valeur, car les cordes courtes, pour les petits $\theta$ et proches de $2\pi$, la tirent vers le bas. Une simulation confirme $1.273$.

\iqwrong{Utiliser la moyenne quadratique. L'espérance du carré de la distance vaut exactement 2, donc sa racine vaut $1.41$, ce qui est trop élevé. La moyenne d'une racine n'est pas la racine d'une moyenne.}
\end{solution}

\begin{solution}{iq:iv1:timed-test-4:6}
La moyenne vaut $\frac94$ tours.

À trois joueurs, il existe 27 issues. Toutes identiques~: 3. Toutes différentes~: 6. Une rejouée a donc probabilité $\frac13$. Sinon, exactement deux symboles apparaissent. Alors un seul gagnant a probabilité $\frac13$ (9 issues), et deux gagnants ont également probabilité $\frac13$ (9 issues).

À deux joueurs, une égalité a probabilité $\frac13$, donc $E_2 = \frac32$.

Pour trois\index{analyse au premier pas},
\[
E_3 = 1 + \frac13E_3 + \frac13E_2 = 1 + \frac13E_3 + \frac12,
\]
ce qui donne $\frac23E_3 = \frac32$, donc $E_3 = \frac94$. Une simulation du jeu renvoie $2.25$.

\iqfollow{À quatre joueurs, quelle est la probabilité d'une partie rejouée~?}{Exactement deux symboles apparaissent dans $3(2^4 - 2) = 42$ des 81 issues, donc une partie rejouée a probabilité $\frac{39}{81} = \frac{13}{27}$, soit presque une chance sur deux. Les grands groupes s'enlisent.}
\end{solution}

\begin{solution}{iq:iv1:timed-test-4:7}
La probabilité vaut $\frac{7}{10}$.

Définissons les états\index{analyse au premier pas} de la série en cours~: $h$ après un pile, $t$ après un face, $u$ après deux faces, et le départ. Chaque symbole désigne la probabilité que le pile gagne depuis cet état~:
\[
h = \frac12 + \frac12t, \qquad t = \frac12h + \frac12u, \qquad u = \frac12h.
\]
Alors $t = \frac12h + \frac14h = \frac34h$, et $h = \frac12 + \frac38h$, ce qui donne $h = \frac45$. Donc $t = \frac35$.

Depuis le départ, le premier lancer mène à $h$ ou à $t$ avec la même probabilité~: $\frac12(\frac45 + \frac35) = \frac{7}{10}$. Une résolution exacte et une simulation concordent.

Ainsi, une cible d'une série plus courte d'un lancer gagne dans $70~\%$ des cas, un peu plus que ce que suggère la probabilité brute de chaque série. Cas limite~: face à quatre faces consécutives, la probabilité de la série de piles serait encore plus élevée.

\iqwrong{Comparer directement $\frac14$ et $\frac18$, ce qui donne $\frac23$. Cela ignore qu'une tentative ratée sur une série se réinitialise différemment, et qu'une série de piles réinitialise la série de faces.}
\end{solution}

\begin{solution}{iq:iv1:timed-test-4:8}
Le nombre espéré vaut $1 + \ln 2 \approx 1.693$.

Prenons les logarithmes. $E_i = -\ln U_i$ suit une loi exponentielle de moyenne 1, et le produit passe sous $\frac12$ lorsque $\sum E_i$ dépasse $\ln 2$.

Les $E_i$ sont donc les intervalles d'un processus de Poisson de taux 1\index{processus de Poisson}. Le nombre d'arrivées dans $[0, \ln 2]$ suit une loi de Poisson de moyenne $\ln 2$. La première arrivée après $\ln 2$ correspond au tirage d'arrêt. Le nombre de tirages vaut donc ce compte de Poisson augmenté de 1, dont la moyenne est $1 + \ln 2$. Une simulation donne $1.693$.

Une seconde méthode~: notons $m(x)$ le nombre espéré de tirages à partir du seuil $x$, de sorte que $m(x) = 1 + \int_x^1m(x/u)\,du$. Cela donne $m(x) = 1 - \ln x$, et en $\frac12$ on retrouve la même valeur.

\iqwrong{Utiliser $\mathbb{E}\prod U_i = 2^{-n}$ et résoudre pour $n = 1$. Il s'agit d'une propriété sur les moyennes, et non sur le premier passage.}

\iqfollow{Et si la cible vaut $e^{-5}$~?}{$1 + 5 = 6$ tirages, par le même raisonnement.}
\end{solution}

\begin{solution}{iq:iv1:timed-test-4:9}
Miser $\frac14$ de sa richesse, pour une croissance logarithmique d'environ $0.059$ par pari.

Une mise de $f$ multiplie la richesse par $1 + 2f$ ou par $1 - f$. Le logarithme espéré vaut
\[
g(f) = \frac12\ln(1 + 2f) + \frac12\ln(1 - f).
\]
En annulant la dérivée~: $\frac{1}{1 + 2f} = \frac{1}{2(1 - f)}$, donc $f = \frac14$\index{critère de Kelly}. La forme générale est l'edge rapporté à la cote~: $(bp - q)/b = (1 - \frac12)/2$.

On obtient alors $g = \frac12\ln 1.5 + \frac12\ln 0.75 = \frac12\ln 1.125 = 0.0589$. Une recherche directe confirme l'optimum.

Porter la mise au double, soit $f = \frac12$, donne $\frac12\ln(2 \times 0.5) = 0$~: aucune croissance du tout. Un effet de levier excessif sur un bon pari ne laisse donc rien, même avec un rendement espéré positif par pari.

\iqwrong{Miser la totalité de sa richesse, parce que le rendement espéré est positif. La première perte met fin au jeu.}
\end{solution}

\begin{solution}{iq:iv1:timed-test-4:10}
L'équipe la plus forte gagne avec une probabilité d'environ $0.710$, et des équipes de même force jouent $5.81$ matchs en moyenne.

Imaginons que les sept matchs soient tous joués. L'équipe remporte la série exactement lorsqu'elle gagne au moins 4 des 7\index{jeu étendu}. On obtient
\[
\sum_{k = 4}^7\binom7k0.6^k0.4^{7 - k} = 0.2903 + 0.2613 + 0.1306 + 0.0280 = 0.710.
\]

Pour des équipes de même force, la série s'achève après $n$ matchs lorsque le vainqueur remporte le match $n$ et 3 des $n - 1$ précédents. En comptant les deux équipes~:
\[
P(n) = 2\binom{n - 1}{3}2^{-n}, \qquad P(4) = \frac18, \quad P(5) = \frac14, \quad P(6) = P(7) = \frac{5}{16}.
\]
La moyenne vaut $\frac{4}{8} + \frac{5}{4} + \frac{30 + 35}{16} = 5.8125$. Une somme exacte sur les matchs confirme les deux résultats.

\iqwrong{Utiliser $\binom{n}{4}$ pour une série qui s'achève au match $n$. Le dernier match doit être gagné, de sorte que seuls les $n - 1$ matchs précédents sont libres. Jouer les sept matchs évite cette erreur, puisque les matchs supplémentaires ne peuvent pas changer le vainqueur.}

\iqwhy{Prolonger un jeu arrêté à une longueur fixe supprime souvent la partie la plus difficile du calcul.}
\end{solution}

\begin{solution}{iq:iv1:timed-test-4:11}
La covariance vaut $\frac{1225}{1296} = (\frac{35}{36})^2 \approx 0.945$.

Comme $Mm$ est le produit des deux dés, $\mathbb{E}Mm = 3.5^2 = 12.25$\index{covariance}.

Le maximum vérifie $P(M \le k) = k^2/36$, ce qui donne $\mathbb{E}M = \frac{161}{36}$. Par la symétrie $x \to 7 - x$, $\mathbb{E}m = 7 - \frac{161}{36} = \frac{91}{36}$.

Donc la covariance vaut $12.25 - \frac{161 \times 91}{1296} = \frac{15\,876 - 14\,651}{1296} = \frac{1225}{1296}$. Une énumération des 36 issues le confirme.

Elle est positive, car les grandes valeurs de l'une vont avec les grandes valeurs de l'autre. Remarquons comment l'identité $Mm = XY$ évite un tableau des 36 produits.

\iqwrong{Supposer que $M$ et $m$ sont non corrélés parce que les dés sont indépendants. Ce sont des fonctions des mêmes deux dés.}
\end{solution}

\begin{solution}{iq:iv1:timed-test-4:12}
La moyenne vaut 50, l'écart type vaut $\sqrt{500} \approx 22.4$, et la probabilité d'un total nul est $e^{-10}$.

Par la loi de l'espérance totale, $\mathbb{E}S = \mathbb{E}N\,\mathbb{E}X = 50$\index{loi de Poisson composée}.

Par la loi de la variance totale,
\[
\operatorname{Var}S = \mathbb{E}N\operatorname{Var}X + \operatorname{Var}N(\mathbb{E}X)^2 = 10 \times 25 + 10 \times 25 = 500.
\]
Pour un comptage de Poisson, cela se simplifie en $\lambda\mathbb{E}X^2 = 10 \times 50$. La moitié de la variance provient du nombre de sinistres, et l'autre moitié de leurs montants. Une simulation le confirme.

Le total n'est nul qu'en l'absence de sinistres, donc la probabilité vaut $e^{-10} \approx 4.5 \times 10^{-5}$.

\iqwrong{Donner $\operatorname{Var}S = \lambda\operatorname{Var}X = 250$. Cela omet l'incertitude sur le nombre de sinistres qui surviennent.}

\iqfollow{Comment le risque relatif change-t-il avec $\lambda$~?}{$\sqrt{\lambda\mathbb{E}X^2}/(\lambda\mathbb{E}X)$ diminue comme $1/\sqrt{\lambda}$. Un portefeuille quatre fois plus grand a la moitié du risque relatif.}
\end{solution}
